% Page 35 du document source
% Transcription provisoire — à relire et corriger si nécessaire.

\begin{center}
    {\Large\itshape Alban au XV\textsuperscript{ème} siècle (suite)}
\end{center}

\vspace{1.5em}

En cette même année 1447, Bertrand Rauly d'Alban, fait hommage au comte
de Castres, pour un fief qu'il tient noblement dans le lieu d'Alban. Voilà
comment cette reconnaissance de possession s'accomplissait, si du moins
nous en croyons un hommage semblable rendu en 1443 à Philippe Jean de
Rabastens, seigneur de Paulin.

\medskip

Le requérant\footnote{L'emploi dans ce contexte mérite vérification.}, à genoux, tête découverte, sans ceinture à la taille,
mettait ses mains entre les mains de celui à qui il rendait hommage :
son suzerain. Le vassal, selon le texte latin, se tenait « genibus flexis,
sine capucis et zona, junctis manibus suis inter manus ipsius domini... ».

\medskip

Un événement plus important, du moins au point de vue ecclésiastique,
c'est regroupement de paroisses ordonné en ce début de siècle par
Monseigneur Robert Dauphin, évêque d'Albi. Négremont se voit dépossédé,
au bénéfice de Notre Dame d'Alban, de son titre d'Archiprêtré (nous
disons aujourd'hui « doyenné »). Le Pape Eugène IV, le 8 janvier 1438,
confirme la décision de l'évêque… Mais le dernier Archiprêtre de
Négremont recevra en compensation une rente viagère. (De nos jours, on
se contente généralement de coiffer le « démissionné » de quelque titre
plus ou moins honorifique… et encore… ne comportant pas de soulagement
en espèces sonnantes et trébuchantes… Autres temps… autres mœurs…).
Quant aux revenus de la paroisse de Négremont, alors 2000 livres, ils
devront assurer le traitement du successeur devenu simple « Vicaire
perpétuel », et payer en partie l'entretien des enfants de chœur de la
cathédrale Sainte Cécile d'Albi.
